Distributed last-level cache (LLC) architectures in modern many-core processors introduce non-uniform latencies for memory access, exploitable as side-channel attacks.
These attacks use physical distance or network congestion to leak sensitive information.
We demonstrate a concrete distance-based attack on an \cpumodel{} CPU, extracting AES secret keys from OpenSSL \opensslver{}.
To counter such attacks, we propose \scheme{}, an efficient architectural defense that creates an oblivious LLC by eliminating observable differences in access time for security-sensitive memory instructions.
We present two core mechanisms: \jscheme{}, using calibrated delay, and \bscheme{}, using router bypass.
A randomized-target variant of \bscheme{}, \rscheme{}, trades determinism for lower average latency.
We implement \scheme{} in the Gem5 simulator and evaluate it on PARSEC, Rodinia, and MLPerf Tiny.
On PARSEC and Rodinia, \jscheme{} incurs 7.24\% overhead on average, while \bscheme{} shows no measurable impact over the insecure baseline.
On MLPerf Tiny, overhead for \bscheme{} ranges from $-30.5\%$ to $+9.1\%$ depending on the benchmark; \rscheme{} lowers this to $-41.5\%$ on average.
